\chapter*{Ход выполнения практики}
\addcontentsline{toc}{chapter}{Ход выполнения практики}

\section*{Задания и отметки о выполнении}

\task{Изучить математические модели теплопередачи в сплошных средах 
и их численную реализацию.}

\task{Разработать расчётный модуль на языке C++ для численного 
решения уравнения теплопроводности Фурье--Кирхгофа на основе 
явной конечно-разностной схемы FTCS.}

\task{Разработать графический интерфейс пользователя на языке Python, 
обеспечивающий настройку параметров расчёта и визуализацию 
результатов моделирования.}

\task{Выполнить численное моделирование нестационарной теплопроводности 
для различных материалов и граничных условий с использованием 
разработанного программного комплекса HeatSim3D.}

\task{Провести верификацию разработанной модели путём сравнения 
результатов численного моделирования с аналитическим решением 
нестационарной задачи теплопроводности.}

\task{Подготовить отчёт по практике в соответствии с требованиями 
нормоконтроля кафедры информационных систем и компьютерного 
моделирования.}

\clearpage
\section*{Отзыв о работе обучающегося}

В ходе выполнения практики Богунь М. В. провёл разработку и 
исследование программного комплекса HeatSim3D для численного 
моделирования нестационарной теплопроводности в трёхмерных 
сплошных средах. Был реализован расчётный модуль на языке C++ 
на основе явной конечно-разностной схемы FTCS, разработан 
графический интерфейс пользователя, а также выполнена верификация 
модели путём сравнения с аналитическим решением.

Студент продемонстрировал умение самостоятельно разрабатывать 
численные алгоритмы, реализовывать их программно и анализировать 
результаты моделирования. При выполнении работы были применены 
современные инструменты разработки программного обеспечения, 
в том числе технология параллельных вычислений OpenMP.

Богунь М. В. проявил высокий уровень самостоятельности и 
ответственности при выполнении всех задач практики и в полной 
мере освоил предусмотренные учебным планом компетенции. 
По результатам практики студент заслуживает оценки «отлично».

\progressapprov